% Fragment partage : inclus par trajets-confiance.tex (dossier autonome)
% et par dossier-conception-alanya.tex (dossier client assemble).
% Ne pas y mettre de preambule ni de \part : c'est l'appelant qui les pose.

\chapter{Le besoin}

\section{Ce qui existe déjà}

Alanya sait partager \textbf{une} position. Depuis le menu des pièces jointes d'une
conversation, on ouvre une carte, on déplace le pin, on envoie. Le destinataire reçoit une
vignette~; un appui l'ouvre dans son application de cartographie.

\begin{constat}[Un partage ponctuel, complet et fonctionnel]
Le sélecteur est \code{lib/\allowbreak screens/\allowbreak chats/\allowbreak location\_\allowbreak picker\_\allowbreak screen.dart}~: carte
OpenStreetMap, pin fixe au centre, et une chaîne de repli soignée --- service désactivé,
permission refusée, \code{getCurrentPosition} trop lent, dernière position connue. Le
message est un message ordinaire de \textbf{type~5}, dont le \code{content} est un JSON
\code{\{lat, lng, name, address\}} produit par
\code{lib/\allowbreak core/\allowbreak utils/\allowbreak location\_payload.dart}, avec un géocodage inverse Nominatim en
meilleur effort.
\end{constat}

Ce qui n'existe pas~: le \emph{mouvement}. Une position envoyée est une photographie. Elle
ne dit pas si la personne avance, si elle s'est arrêtée, ni si elle est arrivée.

\section{Pourquoi le message ne convient pas au temps réel}

L'idée la plus immédiate serait de renvoyer un message de type~5 toutes les trente secondes.
Elle ne tient pas~:

\begin{itemize}[leftmargin=1.4em]
  \item la conversation deviendrait illisible --- des dizaines de vignettes pour un seul
        trajet~;
  \item chaque envoi remonterait la conversation en tête de liste, rallumerait le compteur
        de non-lus et déclencherait une notification~;
  \item l'historique serait impossible à tenir~: on ne saurait pas dire où commence et où
        finit un trajet, ni le supprimer sans supprimer des messages~;
  \item et surtout, rien là-dedans ne \emph{surveille} l'arrivée. Or c'est cela, la demande.
\end{itemize}

\section{Ce qu'on ajoute}

Un objet à part~: le \textbf{trajet}. On le démarre, on annonce où l'on va et quand on
compte arriver, ses proches le suivent, et il se termine par une confirmation --- ou par une
alerte.

\begin{note}[La phrase à retenir]
La promesse n'est pas «~vos proches voient où vous êtes~». Elle est~:
\textbf{«~si vous ne confirmez pas votre arrivée, vos proches seront prévenus~»}.

La carte est un confort. L'échéance est le contrat. C'est une distinction de conception, pas
une nuance de vocabulaire~: elle décide de tout ce qui suit, à commencer par le fait qu'une
perte de signal n'est jamais une alerte.
\end{note}

\section{La promesse, et ce qu'elle n'est pas}

\begin{center}
\begin{tabularx}{\linewidth}{@{}p{4.6cm}X@{}}
\toprule
\thd{Ce n'est pas un traceur} & Il n'existe aucun partage permanent, aucune récurrence,
aucun interrupteur «~toujours actif~». Un trajet est fini, et démarré à la main. \\
\thd{Ce n'est pas un service d'urgence} & Le SOS prévient le cercle de confiance. Il
n'appelle ni la police, ni les secours, et l'écran le dit. \\
\thd{Ce n'est pas une demande} & Personne ne peut réclamer la position de quelqu'un. Seul
le propriétaire démarre, et seul lui arrête. \\
\thd{Ce n'est pas un historique partagé} & Les membres du cercle voient le trajet
\emph{pendant} qu'il a lieu. Ils n'ont accès à aucun trajet passé. \\
\bottomrule
\end{tabularx}
\end{center}

\section{Les règles produit}

\begin{center}
\begin{tabularx}{\linewidth}{@{}p{4.6cm}X@{}}
\toprule
\thd{Qui reçoit} & \textbf{Tout le cercle de confiance}, c'est-à-dire la liste système
\emph{Confiance}, plafonnée à cinq. Pas de sous-sélection au départ. \\
\thd{Un seul à la fois} & Un utilisateur n'a qu'un trajet ouvert. Repartir, c'est un
nouveau trajet. \\
\thd{L'arrivée ne clôt rien} & Le GPS \emph{propose}, l'humain \emph{dispose}. Une
destination atteinte pose la question~; elle n'y répond pas. \\
\thd{Seul le propriétaire clôt} & Un membre peut dire «~j'ai vu~» et appeler. Il ne peut
pas décider que quelqu'un est arrivé. \\
\thd{Dix minutes de grâce} & Entre l'échéance et l'alerte. Avec trois relances, adressées
au propriétaire seul. \\
\thd{Une alerte ne se retire pas} & Elle se \emph{résout}. Le journal garde la trace de ce
qui est parti, et à qui. \\
\thd{Silence n'est pas danger} & Un GPS perdu s'affiche en gris, jamais en rouge. Le rouge
est réservé à l'alerte et au SOS. \\
\thd{Douze heures au maximum} & Passé ce plafond, le trajet demande à être clos. \\
\thd{Navigation} & \textbf{Aucun sixième onglet.} La barre est codée en dur à cinq onglets
(\code{lib/\allowbreak screens/\allowbreak home/\allowbreak glass\_nav\_bar.dart}). \\
\bottomrule
\end{tabularx}
\end{center}

\chapter{Le cercle de confiance}

\section{Pourquoi réutiliser la liste Confiance}

Le volet~2 a créé quatre listes système, dont une~: \emph{Confiance}, seule à porter un
plafond --- cinq membres. Elle attendait un usage~; c'est celui-ci.

\begin{constat}[Le plafond existe déjà en base]
\code{src/\allowbreak utils/\allowbreak defaultContactLists.js} sème la liste \code{kind = 'trust'} avec
\code{member\_limit = 5} et la couleur \code{\#B7791F}, à l'inscription et en rattrapage à
chaque lecture des listes. Le contrôleur refuse le sixième membre en 409
\code{LIST\_MEMBER\_LIMIT}. Rien n'est à créer côté audience.
\end{constat}

Ne pas inventer une seconde notion d'audience est aussi la position déjà tranchée au
volet~2, qui a refusé de fusionner listes personnelles et audiences de diffusion. Ici, on va
dans le même sens~: on \emph{réutilise} la liste, on n'en fabrique pas une autre.

\section{Le cercle est l'audience --- en entier}

Démarrer un trajet prévient les cinq. Il n'y a pas de case à cocher, pas d'écran de
sélection, pas de «~partager à trois sur cinq~».

\begin{note}[Ce que ce choix achète]
Un appui de moins, et une promesse énonçable en une phrase --- «~mon cercle est prévenu~» —
au lieu d'une liste variable dont on ne se souvient plus le lendemain. C'est aussi ce qui
permet de tenir l'objectif de vingt secondes entre l'envie de partir et le départ effectif.

Le corollaire~: on compose son cercle \textbf{à froid}, dans ses listes de contacts, pas au
moment de partir. Cette composition ne notifie personne.
\end{note}

\section{Ce que le membre découvre}

\begin{alerte}[Une règle du volet~2 est mise en défaut --- assumons-le]
Les listes de contacts sont aujourd'hui \textbf{strictement privées}~: «~personne ne sait
dans quelle liste il est rangé, ni même qu'il l'est~»
(\code{contenu/listes-naive.tex}). Le premier trajet partagé \textbf{rompt cette règle}~:
en recevant la carte, un proche apprend qu'il fait partie du cercle de confiance.

On ne peut pas l'éviter --- prévenir quelqu'un sans qu'il sache qu'il a été choisi n'a pas
de sens. On l'assume donc, et l'écran du membre le dit explicitement~:
«~\emph{Awa vous a choisi comme personne de confiance}~». Ce qui reste privé~: les
\emph{autres} listes, et le fait d'être \emph{sorti} du cercle.
\end{alerte}

Le membre voit combien de personnes suivent, jamais qui. Cela rassure --- «~je ne suis pas
seul à veiller~» --- sans exposer le carnet d'adresses de quelqu'un d'autre.

\section{Ce que le membre peut refuser}

Recevoir un trajet n'est pas une obligation. Le membre peut quitter un trajet en cours, et
un réglage général permet de ne plus être destinataire de trajets. Dans les deux cas, le
propriétaire lit «~un membre n'est plus destinataire~», sans motif.

\chapter{Les écrans}

\section{Vue d'ensemble}

Deux chemins partent de la même feuille de départ~: celui où l'on confirme, et celui où le
silence déclenche l'alerte.

\begin{center}
  \imageOuCadre{maquettes/trajets-parcours.png}{\linewidth}%
    {Le chemin nominal (partir, suivre, confirmer, clore) et le chemin d'alerte
     (relances, alerte, prise en charge, résolution)\\
     \texttt{maquettes/trajets-parcours.png}}
\end{center}

\begin{note}[Maquette interactive]
\href{https://claude.ai/code/artifact/fd08532c-8cdb-454d-9af7-8c14a43b79b3}%
{\texttt{claude.ai/code/artifact/fd08532c}}\\
Source versionnée~: \code{docs/\allowbreak architecture/\allowbreak maquettes/\allowbreak trajets-parcours.html}
\end{note}

Les quatre temps --- \textbf{partir, suivre, confirmer, clore} --- sont rappelés en tête de
chaque écran du propriétaire par une rampe de quatre segments. C'est la seule chose qui ne
change jamais de place.

\section{Partir}

\begin{center}
  \imageOuCadre{maquettes/trajets-depart.png}{\linewidth}%
    {Cercle vide, cercle prêt, explication de la permission, composition, puis le contrat\\
     \texttt{maquettes/trajets-depart.png}}
\end{center}

\begin{note}[Maquette interactive]
\href{https://claude.ai/code/artifact/39bf86b6-9506-4aba-937a-6ae051260766}%
{\texttt{claude.ai/code/artifact/39bf86b6}}\\
Source versionnée~: \code{docs/\allowbreak architecture/\allowbreak maquettes/\allowbreak trajets-depart.html}
\end{note}

Trois lignes suffisent~: le type (taxi ou rendez-vous), le lieu, l'heure. On peut n'en
donner qu'une des deux dernières --- avec les deux, la première atteinte pose la question.

Le dernier écran affiche \textbf{le contrat en toutes lettres}~: «~\emph{Vos cinq proches
verront votre position jusqu'à 21:45. Si vous n'avez pas confirmé à 21:55, ils seront
prévenus avec votre dernière position.}~» C'est le seul élément de tout le parcours qui
garantit que la personne a compris ce qu'elle déclenche. Les deux heures y sont écrites, pas
déduites.

Deux écrans ne se voient qu'une fois~: la composition du cercle, quand il est vide, et
l'explication de la permission de localisation --- posée \emph{avant} que le système ne la
demande, parce que c'est le seul moyen d'obtenir «~Toujours autoriser~».

\section{Suivre}

\begin{center}
  \imageOuCadre{maquettes/trajets-en-cours.png}{\linewidth}%
    {Trajet actif, échéance atteinte, signal perdu, bandeau persistant, second appareil\\
     \texttt{maquettes/trajets-en-cours.png}}
\end{center}

\begin{note}[Maquette interactive]
\href{https://claude.ai/code/artifact/32f2cae7-c8db-4f8c-b85a-f3eb43e5d644}%
{\texttt{claude.ai/code/artifact/32f2cae7}}\\
Source versionnée~: \code{docs/\allowbreak architecture/\allowbreak maquettes/\allowbreak trajets-en-cours.html}
\end{note}

Pendant tout le trajet, un \textbf{bandeau} se pose au-dessus de la barre de navigation. Il
est visible sur les cinq onglets, affiche l'heure d'arrivée, et s'ouvre d'un appui.

\begin{constat}[Le bandeau tient dans un espace déjà réservé]
\code{lib/\allowbreak screens/\allowbreak home/\allowbreak glass\_nav\_bar.dart} construit sa rangée par
\code{List.generate(5, …)} et réserve \code{kGlassNavBarSpace = 96} sous le contenu. Le
bandeau s'insère là. Aucun sixième onglet, aucune refonte de navigation.
\end{constat}

«~Arrêter le partage~» est toujours à un appui, y compris depuis la notification système.
Ce n'est pas une commodité~: si arrêter demandait d'ouvrir l'application et de naviguer,
arrêter deviendrait coûteux --- donc reprochable par quelqu'un qui regarde par-dessus
l'épaule.

\section{Côté membre}

\begin{center}
  \imageOuCadre{maquettes/trajets-membre.png}{\linewidth}%
    {La carte de trajet dans la conversation, et ses huit états\\
     \texttt{maquettes/trajets-membre.png}}
\end{center}

\begin{note}[Maquette interactive]
\href{https://claude.ai/code/artifact/065db5f1-a69a-4ce1-af06-efa4ef13b29a}%
{\texttt{claude.ai/code/artifact/065db5f1}}\\
Source versionnée~: \code{docs/\allowbreak architecture/\allowbreak maquettes/\allowbreak trajets-membre.html}
\end{note}

Le trajet arrive chez le proche \textbf{dans la conversation}, là où il parle déjà avec la
personne. C'est la seule surface qu'on est certain qu'il verra~: elle reçoit déjà la
notification, elle survit au redémarrage, elle reste dans l'archive.

Un \textbf{seul message} par trajet, qui change d'état sur place --- il ne remonte pas la
conversation à chaque position. Il reste compact~: avatar, ce qui se passe, l'heure
d'arrivée, et un bouton «~Suivre~». La vraie carte est plein écran, à un appui.

Le membre a exactement deux gestes~: dire qu'il a vu, et appeler.

\section{Décider et alerter}

\begin{center}
  \imageOuCadre{maquettes/trajets-alerte.png}{\linewidth}%
    {Armement du SOS, décompte annulable, mode discret --- puis la réception de l'alerte\\
     \texttt{maquettes/trajets-alerte.png}}
\end{center}

\begin{note}[Maquette interactive]
\href{https://claude.ai/code/artifact/91c2ad6f-cc73-4233-9d55-4b8b905bca8d}%
{\texttt{claude.ai/code/artifact/91c2ad6f}}\\
Source versionnée~: \code{docs/\allowbreak architecture/\allowbreak maquettes/\allowbreak trajets-alerte.html}
\end{note}

Le \textbf{SOS} est accessible sans trajet en cours --- il en crée un, sans destination ni
heure, immédiatement en alerte. Il demande deux gestes pour partir~: un maintien de trois
secondes, puis un décompte annulable de cinq. Après l'envoi, l'écran redevient neutre~:
aucun son, aucune vibration. C'est le moment le plus dangereux, et il ne doit rien laisser
paraître.

Une fausse alerte se \emph{corrige} en un appui, en moins d'une minute. Elle ne s'efface pas.

\section{Après le trajet}

\begin{center}
  \imageOuCadre{maquettes/trajets-fin-historique.png}{\linewidth}%
    {Récapitulatif, historique, historique vide, et l'écran de confidentialité\\
     \texttt{maquettes/trajets-fin-historique.png}}
\end{center}

\begin{note}[Maquette interactive]
\href{https://claude.ai/code/artifact/048b398f-985e-45dd-b4e5-0d34a7ae5a69}%
{\texttt{claude.ai/code/artifact/048b398f}}\\
Source versionnée~: \code{docs/\allowbreak architecture/\allowbreak maquettes/\allowbreak trajets-fin-historique.html}
\end{note}

Le récapitulatif rejoue la frise du trajet --- y compris ce qui a raté, comme les quatre
minutes sans signal. L'historique liste les trajets passés, avec leur issue lisible à la
pastille avant même de lire la ligne.

\section{Quand ça se dégrade}

\begin{center}
  \imageOuCadre{maquettes/trajets-erreurs-admin.png}{\linewidth}%
    {Permission refusée, batterie faible, hors ligne, cercle vidé --- et la vue
     d'administration, réduite à des compteurs\\
     \texttt{maquettes/trajets-erreurs-admin.png}}
\end{center}

\begin{note}[Maquette interactive]
\href{https://claude.ai/code/artifact/556526a1-c27a-4310-9435-158158a106ee}%
{\texttt{claude.ai/code/artifact/556526a1}}\\
Source versionnée~: \code{docs/\allowbreak architecture/\allowbreak maquettes/\allowbreak trajets-erreurs-admin.html}
\end{note}

Aucun de ces quatre écrans n'est rouge. Chacun rappelle que \textbf{l'échéance tient} --- ce
qui est la seule chose que l'utilisateur ait besoin de savoir à ce moment-là, et ce qui
distingue une dégradation d'une panne.

\section{Points d'entrée}

\begin{enumerate}[leftmargin=1.6em]
  \item Le profil, entrée «~Trajets de confiance~», à côté de «~Listes de contacts~»
        (\code{lib/\allowbreak screens/\allowbreak profile/\allowbreak profile\_screen.dart}). C'est l'entrée canonique~:
        démarrer, voir le trajet en cours, relire l'historique.
  \item Le menu des pièces jointes d'une conversation, où «~Position~» devient un choix à
        deux entrées~: \emph{position actuelle} --- le type~5 existant --- ou
        \emph{trajet de confiance}. C'est là que vit déjà le réflexe «~je partage où je
        suis~».
  \item Le bandeau persistant, tant qu'un trajet est ouvert.
\end{enumerate}

\chapter{Les règles de vie d'un trajet}

\section{L'échéance, la grâce, et les trois relances}

\begin{center}
\begin{tabularx}{\linewidth}{@{}p{3.0cm}p{4.4cm}X@{}}
\toprule
\thd{Moment} & \thd{Qui est sollicité} & \thd{Ce qui se passe} \\
\midrule
$T-5$ min          & Le propriétaire, en silence & «~Vous arrivez bientôt~?~» \\
$T$                & Le propriétaire, notification prioritaire &
                     Le trajet passe «~à confirmer~» \\
$T+3$, $T+7$ min   & Le propriétaire, relances & Décompte visible \\
$T+10$ min         & \textbf{Le cercle} & Alerte, avec la dernière position connue \\
\bottomrule
\end{tabularx}
\end{center}

\begin{note}[Pourquoi dix minutes, et non cinq]
En deçà de cinq minutes, chaque embouteillage devient une alerte --- et le cercle apprend à
ne plus les ouvrir. C'est le véritable mode de défaillance de ce genre de produit~: pas
l'alerte manquée, mais l'alerte qu'on ignore. Au-delà d'un quart d'heure, le délai n'a plus
d'utilité dans un incident réel.
\end{note}

Prolonger tient en deux appuis, sans limite de nombre, et le cercle en est informé~:
«~\emph{Awa a prolongé de 15 minutes}~». Cela rassure sans alerter.

\section{Le cercle est figé au départ}

Les cinq destinataires sont photographiés au démarrage. Modifier sa liste Confiance pendant
un trajet ne change pas le trajet en cours~: si l'alerte part, on doit pouvoir dire
exactement qui a été prévenu.

\section{Quand quelqu'un bloque, quand le cercle se vide}

Le blocage l'emporte sur tout. Un membre qui bloque le propriétaire cesse d'être
destinataire --- alertes comprises. Le propriétaire lit «~un membre n'est plus
destinataire~», jamais «~X vous a bloqué~»~: l'application ne révèle pas un blocage.

\begin{alerte}[Un trajet sans destinataire ne continue pas]
S'il ne reste plus personne, le trajet ne doit pas se poursuivre en silence~: il ne serait
plus une fonctionnalité de sécurité, mais un traceur personnel. Le propriétaire reçoit une
invite immédiate~; sans réponse au bout de cinq minutes, le trajet se clôt tout seul.
\end{alerte}

\section{Ce qu'on garde, et combien de temps}

\begin{center}
\begin{tabularx}{\linewidth}{@{}p{4.0cm}p{3.2cm}p{3.2cm}X@{}}
\toprule
 & \thd{Trajet normal} & \thd{Après une alerte} & \thd{Pourquoi} \\
\midrule
La trace détaillée & 24 h après clôture & 30 jours &
Un registre permanent des déplacements de tous les utilisateurs n'est justifié par aucun
besoin. Après un incident, la trace peut servir de preuve. \\
Le trajet lui-même & 12 mois & 12 mois &
L'historique se contente du départ, de l'arrivée, de la durée et de l'issue. \\
\bottomrule
\end{tabularx}
\end{center}

\begin{note}[Le verrou de trente jours protège la personne, pas le service]
Un trajet clos par une alerte ne peut pas être supprimé avant trente jours. C'est
délibéré~: personne ne doit pouvoir contraindre un proche à effacer la trace d'un incident
dans la minute. L'écran de confidentialité l'explique, pour que la règle ne se découvre pas
par un refus.
\end{note}

\section{Ce que l'administration voit --- et ne voit pas}

Une seule page~: des \textbf{compteurs agrégés}. Trajets démarrés, taux de confirmation,
nombre d'alertes, durée médiane. Aucune coordonnée, aucun nom, aucune destination, aucun
trajet individuel.

Deux options ont été examinées puis écartées --- un registre d'incidents pseudonymisé, et un
accès exceptionnel à la dernière position sous double validation. Elles aideraient le
support. Elles sont écartées parce que la posture la plus défendable est celle où la donnée
n'est pas accessible~; et parce qu'à vingt-quatre heures de conservation, il n'y aurait de
toute façon presque rien à consulter.

\section{Ce qui n'est pas dans cette étape}

Déclenchement matériel du SOS par le bouton d'alimentation~; appel automatique des
secours~; partage à quelqu'un qui n'est pas dans le cercle, par lien~; trajets récurrents.
Aucun de ces points ne demanderait de changer le modèle de données.

\chapter{Ce qui pourrait mal tourner}

Une fonctionnalité qui diffuse la position de quelqu'un mérite qu'on écrive ses risques
avant son code. Les quatre suivants sont les plus sérieux, et chacun a sa contre-mesure dans
la conception --- pas dans un réglage.

\section{Le détournement par un proche}

C'est le risque qui peut tuer la fonctionnalité~: un conjoint qui exige un trajet à chaque
sortie. Sept garde-fous, tous structurels~:

\begin{enumerate}[leftmargin=1.6em]
  \item aucune demande de position n'existe --- le propriétaire est toujours l'émetteur~;
  \item tout trajet est fini, et démarré à la main~;
  \item la notification persistante \textbf{nomme les destinataires} et ne peut être
        masquée~;
  \item les membres n'ont \textbf{aucun historique} --- «~montre-moi où tu étais mardi~»
        n'est pas une fonctionnalité~;
  \item l'audience se compose à froid, en silence, hors du moment de partir~;
  \item «~Arrêter le partage~» est toujours à un appui, et un arrêt anticipé n'envoie qu'un
        message neutre~;
  \item un trajet clos par une alerte ne s'efface pas avant trente jours.
\end{enumerate}

\begin{alerte}[Ce sont des atténuations, pas une solution]
Aucune de ces mesures n'empêche de contraindre quelqu'un à démarrer un trajet. Elles rendent
la contrainte moins rentable et moins durable. Il faut l'écrire tel quel --- prétendre le
contraire serait la plus dangereuse des erreurs de conception.
\end{alerte}

\section{La fatigue d'alerte}

Une alerte qu'on n'ouvre plus ne protège personne. D'où les trois relances avant que le
cercle n'entende quoi que ce soit, la prolongation en deux appuis, et une graduation stricte
des notifications~: le démarrage est ordinaire, la clôture est silencieuse, seules l'alerte
et le SOS sont bruyants.

\section{La fausse alerte}

Un SOS parti d'une poche détruirait la confiance du cercle plus sûrement qu'une alerte
manquée. Maintien de trois secondes, décompte annulable de cinq, bouton absent de l'écran
d'accueil, et correction poussée au cercle en moins d'une minute.

\section{La batterie et la précision}

Le suivi n'est pas continu~: il se déclenche à la distance parcourue, avec un battement
régulier qui distingue \emph{immobile} de \emph{traceur éteint}. Sous quinze pour cent, il
ralentit et le cercle en est informé.

Quant à la précision, elle est dite plutôt que masquée~: en ville, on est souvent à
soixante mètres près, et l'écran l'affiche. C'est aussi pourquoi l'arrivée se détecte dans
un rayon de cent cinquante mètres, tenu pendant une minute, à faible vitesse --- passer
devant sa rue sans s'arrêter est le faux positif le plus banal.
